\documentclass{article}
\usepackage[margin=1in]{geometry}
\usepackage{amsmath, amssymb, amsthm}
\usepackage{enumitem}
\usepackage{mathtools}
\usepackage{cancel}

% colored links
\usepackage{hyperref}
\hypersetup{
    colorlinks=true,
    linkcolor=blue,
    filecolor=magenta,      
    urlcolor=black!20!BurntOrange,
    }

% Inputting Python code
\usepackage[dvipsnames]{xcolor}
\definecolor{textblue}{rgb}{.2,.2,.7}
\definecolor{textred}{rgb}{0.54,0,0}
\definecolor{textgreen}{rgb}{0,0.43,0}
\usepackage{upquote}
\usepackage{listings}
\lstset{
    language=Python, 
    tabsize=4,
    basicstyle={\ttfamily},
    keywordstyle=\color{textblue},
    commentstyle=\color{textgreen},
    stringstyle=\color{textred},
    frame=none,
    columns=fullflexible,
    keepspaces=true,
    showstringspaces=false,
    xleftmargin=-15mm, % manual adjustment, figure out permanent solution
}

% Drawing figures
\usepackage{tikz}
\usetikzlibrary{calc, shapes.symbols}

% Colored Boxes
\usepackage{tcolorbox}
\tcbuselibrary{skins,hooks}
\usetikzlibrary{shadows}
\usepackage{lipsum}

%Images
\usepackage{graphicx}
    \usepackage{subcaption}
    \usepackage{float}

%Formatting and Spacing
\setitemize[1]{noitemsep, parsep = 5pt, topsep = 5pt}
\setenumerate[1]{label = (\alph*), parsep = 1pt, topsep = 5pt}
\setlength\parindent{0pt}
\linespread{1.1}

%Easy Numbering in case we add or remove questions
\newcounter{counter}
\newcommand{\Exercise}{Exercise \stepcounter{counter}\thecounter}

\newcommand{\Cov}{\text{Cov}}
\newcommand{\trace}{\text{trace}}
\newcommand{\Normal}{\text{Normal}}
\newcommand{\Var}[1]{\operatorname{Var}\left(#1\right)}
\newcommand{\E}[1]{\operatorname{E}\left[#1\right]}
\renewcommand{\Pr}[1]{\operatorname{P}\left(#1\right)}
\newcommand{\oline}[1]{\overleftrightarrow{#1}}

\DeclareMathOperator*{\minimize}{minimize}
\DeclareMathOperator*{\maximize}{maximize}
\DeclareMathOperator*{\argmin}{argmin}
\DeclareMathOperator*{\argmax}{argmax}
\DeclarePairedDelimiter\abs{\lvert}{\rvert}%
\DeclarePairedDelimiter\norm{\lVert}{\rVert}%
\DeclarePairedDelimiterX{\inp}[2]{\langle}{\rangle}{#1, #2}

%Custom Title Fields
\newcommand{\hmwkTitle}{Homework 7}
\newcommand{\hmwkDueDate}{Mar.\ 27, 2024}
\newcommand{\hmwkDueTime}{11:59pm EST}
\newcommand{\hmwkClass}{Advanced Algorithms}
\newcommand{\hmwkClassInstructor}{Professor Dana Randall}
\newcommand{\hmwkSection}{Spring 2024}
\newcommand{\hmwkAuthorName}{Sarthak Mohanty}

%Headers and Footers
\usepackage{fancyhdr}
\usepackage{extramarks}
\pagestyle{fancy} 
\lhead{CS 4540}
\chead{\hmwkClass \ (\hmwkClassInstructor)}
\rhead{\hmwkTitle}
\cfoot{\thepage}
\renewcommand\headrulewidth{0.4pt}
\renewcommand\footrulewidth{0.4pt}

\title{
    \vspace{2in}
    \textbf{\hmwkClass:\ \hmwkTitle}\\
    \normalsize\vspace{0.1in}\small{Due\ on\ \hmwkDueDate\ at \hmwkDueTime}\\
    \vspace{0.1in}\large{\textit{\hmwkClassInstructor\ \hmwkSection}} \\
    \vspace{0.8in}
    \begin{minipage}[t]{0.4\columnwidth}
        {\footnotesize\textbf{As stated in the syllabus, unauthorized use of previous semester course materials is strictly prohibited in this course.
        }}
    \vspace{0.5in}
    \end{minipage}
    \vspace{0.8in}
    \author{\textbf{\hmwkAuthorName}}
    \date{}
}

\begin{document}

\maketitle
\pagebreak

\subsection*{\Exercise}
    It is sometimes of interest to compute the Voronoi diagram of a set of sites, but we are only interested in a portion of the final diagram. In this problem, we'll consider how to compute the Voronoi diagram of a set of points in $\mathbb{R}^2$, but restricted to a given line $\ell$. By rotating and translating space, we may assume that $\ell$ is aligned with the x-axis.
    
    \vspace{3mm}
    You are given a sequence of $n$ sites in the plane $P = \{p_1, \ldots, p_n\}$ sorted in increasing order of their x-coordinates (see Fig. \ref*{fig:vrestriction}(a)). Present an algorithm that computes the Voronoi diagram of $P$, but restricted only to the $x$-axis. (We don't care about the portion of the diagram lying above or below the axis.)
    
    \vspace{3mm}
    Observe that the diagram is a sequence of intervals that subdivide the x-axis. The output consists of a sequence of (at most $n-1$) endpoints of the segments $\{x_1, \ldots, x_m\}$, and each edge is labeled with the index of the associated site corresponding to this interval (see Fig. \ref*{fig:vrestriction}(b)). Your algorithm should run in $O(n)$ time. (\textit{Hint:} Start by proving that the left-to-right order of the labels along the x-axis is consistent with the left-to-right order of the sites.)
    \begin{figure}[htbp]
        \centering
        \includegraphics[width=0.7\linewidth]{images/voronoi_line.png}
        \caption{Restriction of a Voronoi diagram to a line.}
        \label{fig:vrestriction}
    \end{figure}

\subsection*{\Exercise}
    \textit{(Erickson)} Let $P$ be any set of points in the plane. A triangulation of $P$ is a planar straight-line graph whose outer face is the complement of the convex hull of $P$, and whose bounded faces are all triangles (see Figure \ref*{fig:triangulations}).
    \begin{figure}[htbp]
        \centering
        \includegraphics[width=0.5\linewidth]{images/triangulation.png}
        \caption{An example of a triangulation}
        \label{fig:triangulations}
    \end{figure}
    There are several different ways to measure the \textit{quality} of a triangulation. The goal of this question is to prove that the Delaunay triangulation of $P$ is the best possible triangulation of $P$, for a few different definitions of ``best".
    
    \vspace{3mm}
    For example, Lecture 17 of \href{https://www.cs.umd.edu/~mount/754/Lects/754lects.pdf}{David Mount's notes}\footnote{I gave this as part of your lecture notes.} contains a proof that among all triangulations of $P$, the Delaunay triangulation of $P$ has the largest minimum angle.
    \begin{enumerate}[label=(\alph*)]
        \item Let $T$ be any triangulation of $P$. Prove that if every interior angle of $T$ is acute, then $T$ is the Delaunay triangulation of $P$.
        \item For any triangle $\Delta$ with vertices $pqr$, define
        \[ \text{Vol}(\Delta) = \text{area}(\Delta) \cdot (\|p\|^2 + \|q\|^2 + \|r\|^2) \] where $\|(a, b)\|^2 = a^2 + b^2$ is the squared Euclidean norm. For any triangulation $T$, let $\text{Vol}(T) = \sum_{\Delta \in T} \text{Vol}(\Delta)$. Prove that among all triangulations of $P$, the Delaunay triangulation of $P$ minimizes $\text{Vol}(T)$. (\textit{Hint:} Why is the function called "Vol"?)
    \end{enumerate}

\subsection*{\Exercise}
    Consider the following algorithm for \textsc{VertexCover}, run \textsc{DFS}, then output the nodes which are not leaves in the \textsc{DFS} tree. Show that the output is indeed a vertex cover, and that this algorithm gives yet another 2-approximation for the minimum vertex cover.
  
\subsection*{\Exercise}
    We claimed in class that a bad approach to approximating the optimal vertex cover is to use a greedy method that does the following: pick the vertex that covers the most yet uncovered vertices, add it to the cover, and repeat. Give an infinite family of examples (i.e., one for each value of $n$) that shows that this method will not achieve an approximation ratio of 2.

\subsection*{\Exercise}
    We saw a version of the load-balancing problem (on the exam) where we had 2 machines.  Now we have $m$ machines $M_1 \ldots M_m$ and a set of $n$ jobs. Each job $j$ has a processing time $t_j$. We seek to assign each job to one of the machines so that the loads placed on all machines are as ``balanced" as possible. That is, if $A(i)$ are the set of jobs assigned to machine $M_i$, then the load for machine $M_i$ is $T_i = \sum_{j \in A(i)} t_j$. We want to minimize the \textit{makespan} $T = \max_i T_i$.  This problem is NP-complete, so instead you should design an approximation algorithm for the problem. What approximation ratio does the greedy algorithm (assign each job to the least loaded machine) yield? Prove your answer.
\end{document}
